\section{Appui stratégique et méthodologique}
\markboth{4. Appui stratégique et méthodologique}{}

\subsection{Proposition de démarche projet}

\subsubsection*{Méthodologie retenue}

La démarche suit \textbf{CRISP-DM}, avec une adaptation : la phase de déploiement
n'est pas un aboutissement mais une contrainte amont. Le besoin « servir un
bâtiment non relevé » a été posé \emph{avant} la modélisation, et c'est lui qui a
disqualifié les variables en fuite. Une démarche où le déploiement arrive en
dernier aurait laissé la fuite passer jusqu'en production.

Le cycle est itératif dans les faits : la découverte de la fuite en phase de
préparation a provoqué un retour en compréhension du métier, et la contrainte
d'hébergement rencontrée en phase de déploiement a provoqué un retour en
conception.

\subsubsection*{Feuille de route}

\begin{table}[H]
\centering
\small
\begin{tabularx}{\linewidth}{@{}L{1.3cm}L{3.4cm}XL{2.4cm}@{}}
\toprule
\textbf{Phase} & \textbf{Objet} & \textbf{Livrable de sortie} & \textbf{Jalon} \\
\midrule
0 & Cadrage & Plan de travail, dépôt initialisé, règle de traçabilité & Périmètre arrêté \\
1 & Reprise et audit des données & Jeu nettoyé, journal d'écarts ouvert & Données rejouables \\
2 & Préparation partagée & \code{src/features.py}, analyse bivariée & Préparation figée \\
3 & Modélisation et arbitrages & Banc d'essai, conformalisation, artefacts & Modèles exportés \\
4 & Application & Interface, trois cas d'usage & Application locale \\
5 & Industrialisation & Tests, CI/CD, déploiement & \textbf{Application en ligne} \\
6 & Restitution & Rapport, portfolio & Soutenance \\
\bottomrule
\end{tabularx}
\caption{Feuille de route par phases.}
\end{table}

\subsubsection*{Responsabilités}

Le projet est conduit par une personne seule, ce qui n'annule pas la question des
rôles mais la déplace : chaque rôle doit être \emph{explicitement endossé} à un
moment donné, faute de quoi il est oublié.

\begin{table}[H]
\centering
\small
\begin{tabularx}{\linewidth}{@{}L{3.4cm}X@{}}
\toprule
\textbf{Rôle endossé} & \textbf{Ce qu'il a effectivement produit} \\
\midrule
Analyste métier & Reconstitution des besoins, matrice de priorisation, périmètre \\
Auditeur & Grille d'évaluation de l'existant, mise en évidence de la fuite \\
Data scientist & Préparation, banc d'essai, conformalisation \\
Ingénieur MLOps & Tests, CI/CD, versionnement des artefacts, déploiement \\
Concepteur d'interface & Hiérarchie visuelle, formulation métier de l'incertitude \\
\bottomrule
\end{tabularx}
\caption{Rôles endossés au cours du projet.}
\end{table}

\subsection{Aide à la prise de décision}

\subsubsection*{Risques et opportunités}

\begin{table}[H]
\centering
\small
\begin{tabularx}{\linewidth}{@{}L{4.4cm}L{1.6cm}X@{}}
\toprule
\textbf{Risque} & \textbf{Gravité} & \textbf{Réponse apportée} \\
\midrule
Estimation prise pour une mesure & \textbf{Élevée} & Fourchette systématique, usage écrit dans l'interface \\
Dérive de versions entre entraînement et service & \textbf{Élevée} & Versions épinglées et vérifiées par un test dédié \\
Changement des conditions de l'hébergeur & Moyenne & Déjà survenu deux fois ; le code reste portable \\
Bâtiment hors domaine soumis à l'outil & Moyenne & Avertissement non bloquant au-delà du domaine observé \\
Réentraînement dégradant le modèle & Moyenne & Seuils de non-régression bloquants en intégration continue \\
Abandon de maintenance & Faible & Aucune brique à surveiller en continu, par conception \\
\midrule
\textbf{Opportunité} & & \textbf{Condition de réalisation} \\
\midrule
API réutilisable par un tiers & & Acquise, générée par le cadre applicatif \\
Extension à d'autres millésimes & & Le pipeline est rejouable ; coût marginal faible \\
Extension à d'autres villes & & \textbf{Non} : romprait l'hypothèse d'échangeabilité \\
\bottomrule
\end{tabularx}
\caption{Synthèse des risques et opportunités.}
\end{table}

\subsubsection*{Scénarios budgétaires}

Les coûts d'infrastructure sont documentés ci-dessous. Le coût de mise en œuvre
est exprimé en jours-homme ; sa valorisation dépend du taux journalier retenu et
n'est donc pas chiffrée ici.

\begin{table}[H]
\centering
\small
\begin{tabularx}{\linewidth}{@{}L{3.2cm}L{2.4cm}XL{2.2cm}@{}}
\toprule
\textbf{Scénario} & \textbf{Coût récurrent} & \textbf{Ce qu'il apporte} & \textbf{Décision} \\
\midrule
Hébergement gratuit & \textbf{0\,\$} & Application publique, mise en veille automatique après inactivité & \textbf{Retenu} \\
\addlinespace
Abonnement PRO & 9\,\$ / mois & Débloque le conteneur Docker et le matériel CPU dédié & Écarté : n'achète qu'une différence d'outillage \\
\addlinespace
Matériel CPU dédié & $\approx$ 22\,\$ / mois & Supprime le temps de démarrage à froid & Écarté : voir ci-dessous \\
\bottomrule
\end{tabularx}
\caption{Scénarios budgétaires d'infrastructure.}
\end{table}

\begin{constat}{Un arbitrage éclairé par la mesure d'empreinte}
Le troisième scénario a été écarté sur un argument qui n'est pas seulement
financier. Un conteneur maintenu éveillé consomme, à 10~W de puissance moyenne,
environ \textbf{88~kWh par an} — soit \textbf{douze mille fois} l'énergie
consommée par la totalité du banc d'essai. Pour un modèle réentraîné une fois par
an et un usage épisodique, payer pour supprimer un temps de démarrage revient à
financer une consommation continue sans contrepartie d'usage.
\end{constat}

\subsubsection*{Indicateurs de succès}

\begin{table}[H]
\centering
\small
\begin{tabularx}{\linewidth}{@{}L{4.8cm}L{3.2cm}X@{}}
\toprule
\textbf{Indicateur} & \textbf{Cible} & \textbf{Nature} \\
\midrule
Erreur relative médiane hors bloc & $\leq 40\,\%$ & Technique, bloquant en CI \\
Couverture observée des fourchettes & $\geq 75\,\%$ & Technique, bloquant en CI \\
Bornes ordonnées sur tout le parc & 100\,\% & Technique, bloquant en CI \\
Suite de tests au vert avant déploiement & 100\,\% & Technique, bloquant \\
\midrule
Part des émissions couverte par les 20\,\% premiers audits & À mesurer & Métier \\
Taux de confirmation d'une priorisation par l'audit réel & À mesurer & Métier \\
Nombre de bâtiments pré-qualifiés sans déplacement & À mesurer & Métier \\
\bottomrule
\end{tabularx}
\caption{Indicateurs de succès techniques et métier.}
\end{table}

Les indicateurs métier sont marqués « à mesurer » et non renseignés : ils
supposent un retour d'usage qui n'existe pas encore. Les afficher avec des
valeurs inventées les rendrait inutiles.

\subsubsection*{Impacts pris en compte}

\begin{description}[leftmargin=0em,style=nextline,itemsep=0.35em]
  \item[Éthique et réputationnel] Les données sont publiques et nominatives à
        l'échelle du bâtiment. Publier qu'un bâtiment identifié est « le plus
        consommateur de sa catégorie » sur la foi d'une estimation fausse d'un
        facteur~2 serait dommageable pour son propriétaire. C'est la raison
        première pour laquelle la fourchette est affichée au même niveau que
        l'estimation, et non en note de bas de page.
  \item[Légal] Les estimations sont opposables dans un contexte réglementaire.
        L'interface indique explicitement qu'elles ne peuvent servir ni à
        facturer, ni à contractualiser, ni à dimensionner un investissement.
  \item[Données personnelles] Aucune donnée à caractère personnel n'est traitée :
        le jeu décrit des bâtiments, pas des occupants. Le RGPD n'est pas
        applicable au traitement.
  \item[Environnemental] L'empreinte de l'entraînement a été mesurée plutôt
        qu'estimée, et cette mesure a servi un arbitrage réel
        (section~\ref{sec:carbone}).
  \item[Organisationnel] La solution n'ajoute aucune brique à surveiller en
        continu, ce qui est une décision de conception assumée pour un service
        disposant d'un seul utilisateur métier et d'aucune équipe d'exploitation.
\end{description}

\newpage
